%! @@TITLE@@ — Unit @@UNITINT@@, Lesson @@LESSONID@@ — Slides (Beamer)
\documentclass[aspectratio=169, 11pt]{beamer}
\usepackage{@@PREFIX@@-beamer}   % provides \forestheader{} and \sectionlabel[color]{}

\begin{document}

% TITLE SLIDE (CourseName is not defined in beamer, so the name is literal)
{
\setbeamercolor{background canvas}{bg=forest}
\begin{frame}[plain]
  \vspace{2.8cm}\hspace{0.55cm}%
  \begin{minipage}{0.9\textwidth}
    {\color{goldacc}\bfseries\small LESSON @@LESSONID@@}\\[0.5cm]
    {\color{white}\bfseries\LARGE @@TITLE@@}\\[1.2cm]
    {\color{forestmid}\small @@COURSENAME@@~$\cdot$~Unit @@UNITINT@@}
  \end{minipage}
\end{frame}
}

% The deck follows the GRADUAL-RELEASE flow (~11 frames):
%   learning targets → warm-up → hook → notes 1–4 → group-activity launch → debrief → close.
% The vocabulary is on the slides from the notes frames onward; there is no spoiler rule.

% LEARNING TARGETS
\begin{frame}[t, plain]
  \forestheader{Today's targets}
  \sectionlabel{What you will be able to do}
  \begin{itemize}\setlength{\itemsep}{6pt}
    \item TODO: target, naming the vocabulary.
  \end{itemize}
  \vspace{4pt}
  \begin{block}{How today runs}
    Warm-up (5) $\to$ notes together (15) $\to$ group activity (25) $\to$ debrief (10) $\to$ close (5).
  \end{block}
\end{frame}

% WARM-UP
\begin{frame}[t, plain]
  \forestheader{Warm-Up: three things you already know}
  \sectionlabel{5 minutes --- alone, then check with a neighbour}
  % TODO: the warm-up items.
\end{frame}

% HOOK
\begin{frame}[t, plain]
  \forestheader{Hook: TODO}
  \sectionlabel[goldacc]{TODO the one-line idea}
  % TODO: the context, the quick questions WITH their answers, and a \begin{block} landing the idea.
\end{frame}

% NOTES 1–4 (repeat this frame four times, one per numbered notes section)
\begin{frame}[t, plain]
  \forestheader{Notes N: TODO}
  \sectionlabel{TODO}
  % TODO: the definition/procedure and the worked example, matching the packet's section exactly.
  % Section 3 is usually the target-misconception one — give it a \begin{block} with the case
  % where the two answers disagree.
\end{frame}

% GROUP ACTIVITY LAUNCH
\begin{frame}[t, plain]
  \forestheader{Group Activity: \emph{TODO Title}}
  \sectionlabel[goldacc]{Groups of 3--4 $\cdot$ 25 minutes}
  % TODO: the four parts as a numbered list, plus a figure if the activity has one.
  \vspace{2pt}
  \textbf{The rule:} for every answer, be ready to show me \emph{where you see it}.
\end{frame}

% DEBRIEF
\begin{frame}[t, plain]
  \forestheader{Debrief: the four things to take away}
  \sectionlabel[redacc]{10 minutes --- correct your own page as we go}
  \begin{enumerate}\setlength{\itemsep}{5pt}
    \item TODO: the fully labelled part-1 answer.
    \item TODO: the crux.
    \item TODO: the comparison and what it licenses.
    \item TODO: the model's concept check.
  \end{enumerate}
\end{frame}

% CLOSE
\begin{frame}[t, plain]
  \forestheader{Close}
  \sectionlabel{What changed today}
  TODO: one or two sentences --- what they could do walking in, what they can do walking out.
  \vspace{5pt}
  \begin{block}{Homework}
    The \textbf{Homework} page in your packet --- due next class. TODO: name what it covers.
  \end{block}
  \vspace{4pt}
  \textbf{Next:} TODO one-line preview.
\end{frame}

\end{document}
